\problema{Decode String}{394}{src/06-sliding-window-stack/394-decode-string.cpp}{M}

Dada una cadena codificada \texttt{s} con el patrón \texttt{k[cadena]},
donde \texttt{k} es un entero positivo, decodifícala devolviendo la cadena
expandida, repitiendo \texttt{k} veces lo que está dentro de los
corchetes. El patrón puede anidarse (corchetes dentro de corchetes) y
también puede haber texto normal fuera de cualquier corchete.

\paragraph{La idea, con un ejemplo de la vida real.} Imagina que estás
armando cajas dentro de cajas, con una nota pegada afuera de cada una que
dice ``repetir el contenido tantas veces''. Cuando abres una caja nueva
(un \texttt{[}), guardas en un lugar seguro lo que llevabas armado hasta
ese momento y cuántas veces hay que repetir lo que va adentro, y empiezas a
armar el contenido de esa caja desde cero. Cuando cierras la caja (un
\texttt{]}), repites lo que armaste adentro la cantidad de veces indicada,
y lo pegas de vuelta a lo que tenías guardado antes de abrirla. Esa ``pila
de cajas pendientes'' es literalmente una \textbf{pila} (stack): cada caja
que abres se apila encima de la anterior, y se va desapilando en orden
inverso conforme se cierran.

\begin{center}
\ttfamily
\begin{tabular}{l}
3[ a 2[ c ] ]\\
"c" repetida 2 veces $\to$ "cc"\\
"a"+"cc" $=$ "acc", repetida 3 veces $\to$ "accaccacc"\\
\end{tabular}
\end{center}

\paragraph{Cómo lo resuelve el código, paso a paso.} Mantenemos una pila
donde cada elemento es un par: la cadena que llevábamos construida
\emph{antes} de entrar a ese nivel de corchetes, y el número de
repeticiones pendiente para ese nivel. También llevamos dos variables de
trabajo: \texttt{actual} (lo que se está construyendo en el nivel presente)
y \texttt{numero} (los dígitos acumulados antes de encontrar un \texttt{[}).
\begin{enumerate}
  \item Si el carácter es un dígito, lo acumulamos en \texttt{numero}
        (multiplicando por diez lo que ya teníamos, para soportar números
        de más de un dígito).
  \item Si es \texttt{[}, empujamos a la pila el par
        \texttt{(actual, numero)}, y reiniciamos \texttt{actual} y
        \texttt{numero} para empezar el nuevo nivel limpio.
  \item Si es \texttt{]}, sacamos de la pila el par que guardamos al
        entrar a este nivel, y reconstruimos \texttt{actual} como la
        cadena anterior más la cadena de este nivel repetida la cantidad
        de veces indicada.
  \item Si es cualquier otro carácter (una letra), simplemente la
        agregamos a \texttt{actual}.
\end{enumerate}
Al terminar de recorrer \texttt{s}, \texttt{actual} contiene la cadena
completamente decodificada.

\lstinputlisting{src/06-sliding-window-stack/394-decode-string.cpp}

\complejidad{$O(n \cdot m)$}{$O(n)$}

\paragraph{¿Qué significa esa complejidad?} \texttt{n} es la longitud de la
cadena codificada y \texttt{m} es la longitud de la cadena final ya
expandida (que puede ser mucho más grande que \texttt{n} si hay muchas
repeticiones anidadas). El tiempo depende de cuánto texto terminamos
copiando al reconstruir cada nivel. El espacio es proporcional a
\texttt{n} porque la pila, en el peor caso, tiene tantos niveles como
corchetes anidados haya.

\paragraph{Consejos para la entrevista (Amazon).} Este problema es un buen
ejemplo de simular recursión con una pila explícita: cada \texttt{[}
equivale a ``entrar a una llamada'' y cada \texttt{]} a ``regresar de
ella''. Vale la pena mencionar que también se puede resolver con
recursión real, pero la versión con pila evita el riesgo de desbordar la
pila de llamadas si hay muchísimos niveles de anidamiento. Un detalle
fácil de pasar por alto: el número antes de un corchete puede tener más de
un dígito (\texttt{"10[a]"}), así que hay que acumular dígitos, no
asumir que siempre es uno solo. Casos límite a mencionar: cadena vacía,
cadena sin ningún corchete (se devuelve tal cual), y anidamiento profundo.
